
\documentclass{article}
\usepackage{colortbl}
\usepackage{makecell}
\usepackage{multirow}
\usepackage{supertabular}

\begin{document}

\newcounter{utterance}

\centering \large Interaction Transcript for game `dond', experiment `coop\_en', episode 10 with reflexion{-}Qwen3.5{-}27B.
\vspace{24pt}

{ \footnotesize  \setcounter{utterance}{1}
\setlength{\tabcolsep}{0pt}
\begin{supertabular}{c@{$\;$}|p{.15\linewidth}@{}p{.15\linewidth}p{.15\linewidth}p{.15\linewidth}p{.15\linewidth}p{.15\linewidth}}
    \# & $\;$A & \multicolumn{4}{c}{Game Master} & $\;\:$B\\
    \hline

    \theutterance \stepcounter{utterance}  
    & & \multicolumn{4}{p{0.6\linewidth}}{
        \cellcolor[rgb]{0.9,0.9,0.9}{
            \makecell[{{p{\linewidth}}}]{
                \texttt{\tiny{[P1$\langle$GM]}}
                \texttt{You have attempted this same task before and did not succeed. These are the plans you wrote after those attempts.} \\
\texttt{Trial \#1: New plan: I failed because I assumed the other player would accept my initial verbal offer and submit a complementary proposal, but they submitted their own secret proposal first, likely claiming items I didn't explicitly secure or creating a conflict. My strategy of making a unilateral verbal offer without confirming the other player's specific valuation or securing a mutual agreement before the deadline was flawed. Next time, I will explicitly ask the other player for their secret value function or their specific item preferences before making a proposal. I will then negotiate a specific split (e.g., "I take the glass, you take the wires and chalks") and wait for their explicit confirmation ("Yes, that works") before submitting my secret proposal. If they submit a proposal first, I must immediately analyze it against the total inventory to ensure my counter{-}proposal is strictly complementary, rather than assuming my previous verbal offer was binding.} \\
\\ 
\\ 
\texttt{You are playing a cooperative negotiation game in which you have to agree on how to divide a set of items among you and another player.} \\
\\ 
\texttt{Rules:} \\
\texttt{(a) You and the other player are given a set of items. Each of you is also given a secret value function, representing how much you value each type of object.} \\
\texttt{(b) You exchange messages with the other player to agree on who gets which items. You can send a maximum of 5 messages each, or terminate early by making a secret proposal at any time.} \\
\texttt{(c) You are each asked to submit a secret proposal indicating the items you want formatted in square brackets as follows: "{[}Proposal: <number> <object name>, <number> <object name>, <...>{]}"} \\
\texttt{(d) If your proposals are complementary, i.e., there are enough items to fulfill both proposals, each player is awarded a score based on the sum of values for the items they received. Otherwise, both of you get zero points.} \\
\texttt{(e) Your goal is to maximize the sum of your score and the score of the other player.} \\
\\ 
\texttt{Let us start.} \\
\\ 
\texttt{The set of available items is:} \\
\\ 
\texttt{1 glass, 2 wires, 2 chalks} \\
\\ 
\texttt{Your secret value function is:} \\
\\ 
\texttt{glass: 8, wire: 1, chalk: 0} \\
\\ 
\texttt{IMPORTANT: Your messages, unless it is the secret proposal, are directly transmitted to the other player, so do not include any response to the rules or text announcing your message. To make a secret proposal, use the indicated format. Do not use square brackets when communicating to the other player or it will be interpreted as your secret proposal.} \\
            }
        }
    }
    & \\ \\

    \theutterance \stepcounter{utterance}  
    & \multicolumn{4}{p{0.6\linewidth}}{
        \cellcolor[rgb]{0.9,0.9,0.9}{
            \makecell[{{p{\linewidth}}}]{
                \texttt{\tiny{[P1$\rangle$GM]}}
                \texttt{Hi! I see we have 1 glass, 2 wires, and 2 chalks to divide. My values are: glass (8), wire (1), and chalk (0). Since I value the glass highly and don't care about chalks, I'd love to get the glass. Could you share your secret values or tell me which items you prefer? That way, we can make sure we get the best total score for both of us.} \\
            }
        }
    }
    & & \\ \\

    \theutterance \stepcounter{utterance}  
    & & & \multicolumn{2}{p{0.3\linewidth}}{
        \cellcolor[rgb]{0.9,0.9,0.9}{
            \makecell[{{p{\linewidth}}}]{
                \texttt{\tiny{[GM$|$GM]}}
                \texttt{continue} \\
            }
        }
    }
    & & \\ \\

    \theutterance \stepcounter{utterance}  
    & & & \multicolumn{2}{p{0.3\linewidth}}{
        \cellcolor[rgb]{0.9,0.9,0.9}{
            \makecell[{{p{\linewidth}}}]{
                \texttt{\tiny{[GM$|$GM]}}
                \texttt{Hi! I see we have 1 glass, 2 wires, and 2 chalks to divide. My values are: glass (8), wire (1), and chalk (0). Since I value the glass highly and don't care about chalks, I'd love to get the glass. Could you share your secret values or tell me which items you prefer? That way, we can make sure we get the best total score for both of us.} \\
            }
        }
    }
    & & \\ \\

    \theutterance \stepcounter{utterance}  
    & & \multicolumn{4}{p{0.6\linewidth}}{
        \cellcolor[rgb]{0.9,0.9,0.9}{
            \makecell[{{p{\linewidth}}}]{
                \texttt{\tiny{[GM$\rangle$P2]}}
                \texttt{You have attempted this same task before and did not succeed. These are the plans you wrote after those attempts.} \\
\texttt{Trial \#1: New plan: I failed because I submitted a proposal claiming only the glass (`{[}Proposal: 1 glass{]}`) while the other player's proposal (implied by the negotiation) likely claimed the chalks and wires. Since the rules state that proposals must be *complementary* (covering all items without overlap) to score, leaving the chalks and wires unclaimed in my proposal caused a mismatch, resulting in zero points for both. I mistakenly assumed that agreeing verbally was sufficient or that the other player would fill the gaps, but the game requires my specific proposal to explicitly list *all* items I am claiming. Next time, after agreeing on the split (Me: 2 chalks, 2 wires; You: 1 glass), I will submit a proposal that explicitly lists every item I am taking: `{[}Proposal: 2 chalks, 2 wires{]}`. This ensures the union of our proposals covers the entire set \{1 glass, 2 wires, 2 chalks\} exactly once.} \\
\\ 
\\ 
\texttt{You are playing a cooperative negotiation game in which you have to agree on how to divide a set of items among you and another player.} \\
\\ 
\texttt{Rules:} \\
\texttt{(a) You and the other player are given a set of items. Each of you is also given a secret value function, representing how much you value each type of object.} \\
\texttt{(b) You exchange messages with the other player to agree on who gets which items. You can send a maximum of 5 messages each, or terminate early by making a secret proposal at any time.} \\
\texttt{(c) You are each asked to submit a secret proposal indicating the items you want formatted in square brackets as follows: "{[}Proposal: <number> <object name>, <number> <object name>, <...>{]}"} \\
\texttt{(d) If your proposals are complementary, i.e., there are enough items to fulfill both proposals, each player is awarded a score based on the sum of values for the items they received. Otherwise, both of you get zero points.} \\
\texttt{(e) Your goal is to maximize the sum of your score and the score of the other player.} \\
\\ 
\texttt{Let us start.} \\
\\ 
\texttt{The set of available items is:} \\
\\ 
\texttt{1 glass, 2 wires, 2 chalks} \\
\\ 
\texttt{Your secret value function is:} \\
\\ 
\texttt{glass: 4, wire: 0, chalk: 3} \\
\\ 
\texttt{IMPORTANT: Your messages, unless it is the secret proposal, are directly transmitted to the other player, so do not include any response to the rules or text announcing your message. To make a secret proposal, use the indicated format. Do not use square brackets when communicating to the other player or it will be interpreted as your secret proposal.} \\
\\ 
\\ 
\texttt{Hi! I see we have 1 glass, 2 wires, and 2 chalks to divide. My values are: glass (8), wire (1), and chalk (0). Since I value the glass highly and don't care about chalks, I'd love to get the glass. Could you share your secret values or tell me which items you prefer? That way, we can make sure we get the best total score for both of us.} \\
            }
        }
    }
    & \\ \\

    \theutterance \stepcounter{utterance}  
    & & & \multicolumn{4}{p{0.6\linewidth}}{
        \cellcolor[rgb]{0.9,0.9,0.9}{
            \makecell[{{p{\linewidth}}}]{
                \texttt{\tiny{[GM$\langle$P2]}}
                \texttt{Hi! Thanks for sharing your values. My values are: glass (4), wire (0), and chalk (3).} \\
\\ 
\texttt{Since you value the glass highly (8) and I value it less (4), it makes sense for you to take the **1 glass**.} \\
\\ 
\texttt{I value chalks highly (3 each) and don't care about wires (0), while you don't care about chalks (0) and value wires slightly (1). To maximize our total score, I should take the **2 chalks**, and you should take the **2 wires**.} \\
\\ 
\texttt{This split gives us:} \\
\texttt{{-} You: 1 glass (8) + 2 wires (2) = 10 points} \\
\texttt{{-} Me: 2 chalks (6) = 6 points} \\
\texttt{{-} Total: 16 points} \\
\\ 
\texttt{Does this sound good to you? If so, I will submit my proposal for the 2 chalks, and you can submit yours for the glass and wires.} \\
            }
        }
    }
     \\ \\

    \theutterance \stepcounter{utterance}  
    & & & \multicolumn{2}{p{0.3\linewidth}}{
        \cellcolor[rgb]{0.9,0.9,0.9}{
            \makecell[{{p{\linewidth}}}]{
                \texttt{\tiny{[GM$|$GM]}}
                \texttt{continue} \\
            }
        }
    }
    & & \\ \\

    \theutterance \stepcounter{utterance}  
    & & & \multicolumn{2}{p{0.3\linewidth}}{
        \cellcolor[rgb]{0.9,0.9,0.9}{
            \makecell[{{p{\linewidth}}}]{
                \texttt{\tiny{[GM$|$GM]}}
                \texttt{Hi! Thanks for sharing your values. My values are: glass (4), wire (0), and chalk (3).} \\
\\ 
\texttt{Since you value the glass highly (8) and I value it less (4), it makes sense for you to take the **1 glass**.} \\
\\ 
\texttt{I value chalks highly (3 each) and don't care about wires (0), while you don't care about chalks (0) and value wires slightly (1). To maximize our total score, I should take the **2 chalks**, and you should take the **2 wires**.} \\
\\ 
\texttt{This split gives us:} \\
\texttt{{-} You: 1 glass (8) + 2 wires (2) = 10 points} \\
\texttt{{-} Me: 2 chalks (6) = 6 points} \\
\texttt{{-} Total: 16 points} \\
\\ 
\texttt{Does this sound good to you? If so, I will submit my proposal for the 2 chalks, and you can submit yours for the glass and wires.} \\
            }
        }
    }
    & & \\ \\

    \theutterance \stepcounter{utterance}  
    & & \multicolumn{4}{p{0.6\linewidth}}{
        \cellcolor[rgb]{0.9,0.9,0.9}{
            \makecell[{{p{\linewidth}}}]{
                \texttt{\tiny{[P1$\langle$GM]}}
                \texttt{Hi! Thanks for sharing your values. My values are: glass (4), wire (0), and chalk (3).} \\
\\ 
\texttt{Since you value the glass highly (8) and I value it less (4), it makes sense for you to take the **1 glass**.} \\
\\ 
\texttt{I value chalks highly (3 each) and don't care about wires (0), while you don't care about chalks (0) and value wires slightly (1). To maximize our total score, I should take the **2 chalks**, and you should take the **2 wires**.} \\
\\ 
\texttt{This split gives us:} \\
\texttt{{-} You: 1 glass (8) + 2 wires (2) = 10 points} \\
\texttt{{-} Me: 2 chalks (6) = 6 points} \\
\texttt{{-} Total: 16 points} \\
\\ 
\texttt{Does this sound good to you? If so, I will submit my proposal for the 2 chalks, and you can submit yours for the glass and wires.} \\
            }
        }
    }
    & \\ \\

    \theutterance \stepcounter{utterance}  
    & \multicolumn{4}{p{0.6\linewidth}}{
        \cellcolor[rgb]{0.9,0.9,0.9}{
            \makecell[{{p{\linewidth}}}]{
                \texttt{\tiny{[P1$\rangle$GM]}}
                \texttt{That sounds perfect! The split maximizes our combined score, and we both get the items we value most.} \\
\\ 
\texttt{I will submit my secret proposal for the glass and the wires. Please go ahead and submit yours for the chalks.} \\
\\ 
\texttt{{[}Proposal: 1 glass, 2 wires{]}} \\
            }
        }
    }
    & & \\ \\

    \theutterance \stepcounter{utterance}  
    & & & \multicolumn{2}{p{0.3\linewidth}}{
        \cellcolor[rgb]{0.9,0.9,0.9}{
            \makecell[{{p{\linewidth}}}]{
                \texttt{\tiny{[GM$|$GM]}}
                \texttt{proposal} \\
            }
        }
    }
    & & \\ \\

    \theutterance \stepcounter{utterance}  
    & & & \multicolumn{2}{p{0.3\linewidth}}{
        \cellcolor[rgb]{0.9,0.9,0.9}{
            \makecell[{{p{\linewidth}}}]{
                \texttt{\tiny{[GM$|$GM]}}
                [1, 2, 0]
            }
        }
    }
    & & \\ \\

    \theutterance \stepcounter{utterance}  
    & & \multicolumn{4}{p{0.6\linewidth}}{
        \cellcolor[rgb]{0.9,0.9,0.9}{
            \makecell[{{p{\linewidth}}}]{
                \texttt{\tiny{[GM$\rangle$P2]}}
                \texttt{The other player has submitted a secret proposal. It is now your turn to submit a proposal of your own.} \\
            }
        }
    }
    & \\ \\

    \theutterance \stepcounter{utterance}  
    & & & \multicolumn{4}{p{0.6\linewidth}}{
        \cellcolor[rgb]{0.9,0.9,0.9}{
            \makecell[{{p{\linewidth}}}]{
                \texttt{\tiny{[GM$\langle$P2]}}
                \texttt{{[}Proposal: 2 chalks{]}} \\
            }
        }
    }
     \\ \\

    \theutterance \stepcounter{utterance}  
    & & & \multicolumn{2}{p{0.3\linewidth}}{
        \cellcolor[rgb]{0.9,0.9,0.9}{
            \makecell[{{p{\linewidth}}}]{
                \texttt{\tiny{[GM$|$GM]}}
                \texttt{proposal} \\
            }
        }
    }
    & & \\ \\

    \theutterance \stepcounter{utterance}  
    & & & \multicolumn{2}{p{0.3\linewidth}}{
        \cellcolor[rgb]{0.9,0.9,0.9}{
            \makecell[{{p{\linewidth}}}]{
                \texttt{\tiny{[GM$|$GM]}}
                [0, 0, 2]
            }
        }
    }
    & & \\ \\

    \theutterance \stepcounter{utterance}  
    & & & \multicolumn{2}{p{0.3\linewidth}}{
        \cellcolor[rgb]{0.9,0.9,0.9}{
            \makecell[{{p{\linewidth}}}]{
                \texttt{\tiny{[GM$|$GM]}}
                [[1, 2, 0], [0, 0, 2]]
            }
        }
    }
    & & \\ \\

\end{supertabular}
}

\end{document}
